\section{Cài đặt}
Nhóm cài đặt lại Eclat từ bài báo gốc của Zaki~\cite{zaki2000} bằng Julia, không dùng
thư viện khai thác tập phổ biến (frequent itemset mining, FIM) nào. Chương này trình bày
biểu diễn dữ liệu, thủ tục tìm kiếm, lý do thuật toán cho kết quả đúng, độ phức tạp và
cách nhóm kiểm thử bản cài đặt cơ bản (baseline).

\subsection{Biểu diễn dọc và tidset}
Cơ sở dữ liệu (CSDL) gồm $N$ giao dịch (transaction). Mỗi giao dịch được gán một mã
giao dịch (transaction identifier, tid) là chỉ số $1, 2, \dots, N$ theo thứ tự đọc.
Apriori lưu CSDL theo chiều ngang (horizontal layout): mỗi giao dịch kèm danh sách item.
Eclat dùng biểu diễn dọc (vertical layout): mỗi item $x$ kèm tidset $t(x)$, tức tập mã
các giao dịch chứa $x$. Với itemset $X$, ta có
\begin{equation}
  t(X) = \{\, i \in \{1,\dots,N\} : X \subseteq T_i \,\}, \qquad
  \sigma(X) = |t(X)|,
  \label{eq:support}
\end{equation}
trong đó $T_i$ là giao dịch thứ $i$ và $\sigma(X)$ là độ hỗ trợ (support) của $X$, tức
số giao dịch chứa $X$. Itemset $X$ là phổ biến (frequent) khi
$\sigma(X) \ge \mathit{minsup}$, với $\mathit{minsup}$ là ngưỡng hỗ trợ tối thiểu tính
bằng số giao dịch. Giao dịch rỗng vẫn có tid và vẫn được tính vào $N$, nhưng không xuất
hiện trong tidset nào.

Bản cơ bản lưu mỗi tidset dưới dạng \texttt{Vector\{Int\}} tăng dần, gọi là tid-list.
Hàm \texttt{to\_class} đọc CSDL một lần và thêm tid theo thứ tự tăng dần, nên mọi
tid-list đều đã sắp xếp mà không cần sắp lại.

Ưu điểm chính của biểu diễn dọc là phép giao. Một giao dịch chứa $P \cup \{a, b\}$ khi
và chỉ khi nó chứa cả $P \cup \{a\}$ lẫn $P \cup \{b\}$, nên
\begin{equation}
  t(P \cup \{a, b\}) = t(P \cup \{a\}) \cap t(P \cup \{b\}),
  \label{eq:intersect}
\end{equation}
trong đó $P$ là một itemset và $a, b$ là hai item không thuộc $P$. Độ hỗ trợ của itemset
mới vì vậy được tính từ hai tidset sẵn có, không cần quét lại CSDL.

\subsection{Lớp tương đương theo tiền tố và tìm kiếm Bottom-Up}
Zaki chia không gian tìm kiếm thành các lớp tương đương theo tiền tố (prefix-based
equivalence class)~\cite{zaki2000}. Lớp có tiền tố (prefix) $P$ gồm các phần tử
$(b, t(P \cup \{b\}))$, với mỗi $b$ lớn hơn mọi item của $P$. Hai itemset trong cùng
một lớp chỉ khác nhau ở item cuối, nên phép giao~\eqref{eq:intersect} chỉ hợp lệ giữa
các phần tử cùng lớp.

Hàm \texttt{bottom\_up!} cài đặt thủ tục Bottom-Up của bài báo theo chiều sâu
(depth-first search, DFS). Với mỗi phần tử $a$ của lớp tiền tố $P$, hàm ghi
$P \cup \{a\}$ cùng độ hỗ trợ, giao $t(P \cup \{a\})$ với tidset của mọi phần tử $b$
đứng sau $a$, giữ lại các kết quả phổ biến làm lớp con có tiền tố $P \cup \{a\}$, rồi
duyệt lớp con đó trước khi xét phần tử kế tiếp. Lớp gốc có tiền tố rỗng và gồm các item
phổ biến. Bài báo đếm trước các 2-itemset phổ biến bằng một mảng tam giác trên CSDL
ngang; bản cơ bản bắt đầu phép giao từ các item đơn để một thủ tục xử lý mọi mức.

Bảng~\ref{tab:trace} truy vết thuật toán trên CSDL \texttt{baskets.txt} gồm bốn giao
dịch $T_1 = \{1,2,3\}$, $T_2 = \{1,2\}$, $T_3 = \{1,3\}$, $T_4 = \{2,3\}$ với
$\mathit{minsup} = 2$. Các tidset ban đầu là $t(1) = \{1,2,3\}$, $t(2) = \{1,2,4\}$ và
$t(3) = \{1,3,4\}$. Thuật toán ghi đúng sáu itemset phổ biến, và tập $\{1,2,3\}$ bị tỉa
vì chỉ xuất hiện trong $T_1$.

\begin{table}[t]
\centering
\caption{Truy vết Bottom-Up trên \texttt{baskets.txt} với $\mathit{minsup} = 2$}
\label{tab:trace}
\begin{tabular}{llll}
\hline
Tiền tố & Phép giao & Tidset & Kết quả \\
\hline
$\emptyset$ & -- & $t(1) = \{1,2,3\}$ & ghi $\{1\}$, $\sigma = 3$ \\
$\{1\}$ & $t(1) \cap t(2)$ & $\{1,2\}$ & giữ $2$ \\
$\{1\}$ & $t(1) \cap t(3)$ & $\{1,3\}$ & giữ $3$ \\
$\{1\}$ & -- & $t(12) = \{1,2\}$ & ghi $\{1,2\}$, $\sigma = 2$ \\
$\{1,2\}$ & $t(12) \cap t(13)$ & $\{1\}$ & tỉa, $1 < 2$ \\
$\{1\}$ & -- & $t(13) = \{1,3\}$ & ghi $\{1,3\}$, $\sigma = 2$ \\
$\emptyset$ & -- & $t(2) = \{1,2,4\}$ & ghi $\{2\}$, $\sigma = 3$ \\
$\{2\}$ & $t(2) \cap t(3)$ & $\{1,4\}$ & giữ $3$ \\
$\{2\}$ & -- & $t(23) = \{1,4\}$ & ghi $\{2,3\}$, $\sigma = 2$ \\
$\emptyset$ & -- & $t(3) = \{1,3,4\}$ & ghi $\{3\}$, $\sigma = 3$ \\
\hline
\end{tabular}
\end{table}

\subsection{Tính đúng đắn}
Mọi lớp mà \texttt{bottom\_up!} nhận đều thỏa bốn bất biến (invariant): (i) mỗi tid-list
tăng dần và không lặp; (ii) các phần tử xếp tăng dần theo item và lớn hơn mọi item của
tiền tố; (iii) mỗi phần tử là phổ biến; (iv) tid-list của phần tử $b$ đúng bằng
$t(P \cup \{b\})$. Lớp gốc thỏa cả bốn khi được tạo. Lớp con lấy các phần tử đứng sau
$a$ nên giữ (i) và (ii); phép kiểm tra kích thước trước khi thêm cho (iii); đẳng
thức~\eqref{eq:intersect} cho (iv).

Từ đó suy ra ba tính chất. Thứ nhất, độ hỗ trợ được ghi là chính xác theo (iv)
và~\eqref{eq:support}. Thứ hai, không itemset nào bị ghi hai lần: theo (ii) mọi khóa
$(x_1, \dots, x_k)$ tăng ngặt, và nó chỉ được ghi khi lớp tiền tố
$(x_1, \dots, x_{k-1})$ duyệt tới $x_k$; mỗi lớp chỉ được tạo một lần. Thứ ba, kết quả
đầy đủ nhờ tính đơn điệu giảm (anti-monotone) của độ hỗ trợ: nếu $X \subseteq Y$ thì
$\sigma(X) \ge \sigma(Y)$, nên việc tỉa một itemset không phổ biến không làm mất tập cha
phổ biến nào. Quy nạp theo độ dài cho thấy lớp tiền tố $(x_1, \dots, x_{j-1})$ chứa
mọi item $y > x_{j-1}$ sao cho $(x_1, \dots, x_{j-1}, y)$ phổ biến, nên mọi itemset phổ
biến đều được ghi. Tài liệu \texttt{docs/eclat.md} trong repository trình bày chứng minh
đầy đủ.

\subsection{Độ phức tạp}
Gọi $L$ là tổng số item trên mọi giao dịch, $m$ là số item phân biệt và $C$ là số phép
giao. Dựng biểu diễn dọc tốn $O(L)$ thời gian kỳ vọng cho từ điển, $O(m \log m)$ để sắp
xếp và $O(L)$ bộ nhớ. Một phép giao tốn $O(|t(P \cup \{a\})| + |t(P \cup \{b\})|)$, tối
đa $O(N)$. Một lớp có $k$ phần tử thực hiện $k(k-1)/2$ phép giao, nên tổng thời gian là
\begin{equation}
  O(L + m \log m + C \cdot N),
  \label{eq:cost}
\end{equation}
với $C$ là tổng số phép giao trên mọi lớp được duyệt. Vì tìm kiếm theo chiều sâu, bộ nhớ
chỉ giữ các lớp trên đường đi hiện tại, mỗi mức một lớp gồm tối đa $m$ phần tử, mỗi
phần tử tối đa $N$ tid. Số kết quả có thể tăng theo hàm mũ: một $k$-itemset phổ biến có
$2^k - 1$ tập con khác rỗng, và mọi tập con đó đều phổ biến.

\subsection{Giao diện và định dạng SPMF}
Hai ứng dụng và phần thực nghiệm chỉ gọi bốn hàm công khai: \texttt{read\_transactions}
đọc file SPMF (mỗi dòng một giao dịch, các item cách nhau bởi dấu cách),
\texttt{support\_threshold} đổi $\mathit{minsup}$ thành số giao dịch, \texttt{mine} trả
về từ điển \texttt{SupportMap} từ itemset sang độ hỗ trợ, và \texttt{write\_itemsets}
ghi mỗi itemset một dòng theo dạng \texttt{1 2 \#SUP: 2}. Một $\mathit{minsup}$ thực
trong $(0, 1]$ là tỉ lệ và được làm tròn lên thành $\lceil \mathit{minsup} \cdot N \rceil$
giao dịch; một $\mathit{minsup}$ nguyên là số giao dịch tuyệt đối. Tham số
\texttt{variant} chọn \texttt{:baseline} hoặc \texttt{:optimized}, nên hai phiên bản được
chọn riêng và so sánh trên cùng đầu vào.

Chương trình dòng lệnh \texttt{demo.jl} nhận đường dẫn đầu vào, \texttt{--output},
đúng một trong hai ngưỡng \texttt{--minsup} (tỉ lệ) hoặc \texttt{--min-count} (số giao
dịch) và \texttt{--variant}, ví dụ
\texttt{julia --project=. demo.jl --minsup 0.5 --output out.txt data/toy/baskets.txt}.

\subsection{Kiểm thử và tối ưu hóa}
Bộ kiểm thử tự động (\texttt{julia --project=. test/runtests.jl}) gồm 490 phép kiểm
tra. Với bản cơ bản, bộ kiểm thử so sánh toàn bộ \texttt{SupportMap} trên sáu CSDL nhỏ
có kết quả tính tay (gồm CSDL rỗng, CSDL không có tập phổ biến và giao dịch lặp) và một
CSDL ảnh có giao dịch rỗng. Bộ kiểm thử còn đối chiếu với một bộ đếm vét cạn
(brute force) liệt kê mọi tập con trên 300 CSDL ngẫu nhiên sinh với hạt giống 42, kiểm
tra các ngưỡng biên, một phản ví dụ cho phép giao giữa hai lớp khác tiền tố, thứ tự
khóa và tính ổn định kiểu (type stability). Bộ đếm vét cạn chỉ là công cụ kiểm thử;
phép so sánh với SPMF được trình bày trong chương thực nghiệm.

Bổ sung kỹ thuật tối ưu đã cài đặt kèm bằng chứng đo so với bản cơ bản.
